\chapter{相关技术综述}
\ctexset{section/afterskip=3pt}
\section{Spring Boot}
后端架构选型——本系统选择 Spring Boot。其优势在于：在当前开发生态中具备强大的性能，并通过其专有的“约定优于配置（convention over configuration）”理念大幅提升开发效率。通过引入起步依赖（Starters）以及自动配置机制，改变了整个开发、依赖配置和环境部署繁琐的性质，从而省去在编写业务逻辑之前原本需要经历的中间环节 \cite{Yang2021, Zhao2020}。

Spring Boot 对于实现所需的服务编排功能同样至关重要，用于对科研经费需求进行多维度分析。在该系统中，例如 \texttt{FundStatisticsServiceImpl} 模块可利用其微内核架构，为项目、预算或支出等核心数据实体提供深度且动态的聚合。对于大量经费记录，后端获取模块能够轻松地将底层数据塑造成多维分析对象（例如 \texttt{ExecutionRateVO}）。此外，内置的 Tomcat 容器结合异常处理机制，能够在高并发条件下为数据检索提供强有力的后台支持，从而在统计查询场景中在很大程度上保障决策支持功能的响应速度。
\ctexset{section/afterskip=3pt}
\section{MyBatis-Plus}
在数据持久化方面，系统引入一种名为 MyBatis-Plus 的增强工具。它的设计在保留原有 MyBatis 灵活性的同时，通过提供内置的 BaseMapper 实现了许多自动的 CRUD 操作。例如，针对系统所选择的 3.5.3.1 版本，该版本提供了成熟的 Lambda 条件构建器，开发者可以采用面向对象的设计理念来构建 SQL Dev，并使用更简洁的风格。不仅消除了人为拼写错误的可能性，同时也为持久层代码逻辑提供了更清晰的语义。

MyBatis-Plus 是为在多维度分析中需要进行数据聚合的场景而引入的。传统方式是编写大量 XML 映射文件，而 MP 的 LambdaQueryWrapper 使得在诸如“按类别统计”或“计算总体执行率”等操作变得更为直接。此外，它还利用其自动填充的特性来管理资金流向审计记录，并能够跟踪每一笔支出的完整生命周期。结合生成分布式 ID 的策略，系统能够为 \texttt{FundStatisticsSnapshot} 的快照统计提供稳定的存储支持，并为历史趋势分析建立可靠的数据基础。

\section{Vue 与 ECharts 可视化技术}
Vue 3 是一种常见的前端展示层工具，用于构建响应式界面。本系统采用基于组合式 API（Composition API）的 3.5.12 版本，该版本相较于前代在代码复用与运行时性能方面已被证明更优。本研究选择 ECharts 5.5.1 作为可视化引擎，以便使决策支持仪表板更易于使用。它具备双引擎渲染能力，即可同时使用 Canvas 与 SVG，因此在处理大量财务数据点时，前端仍可实现动态展示 \cite{Xie2019}。

使用 Vue 3 + ECharts 可以将静态图表转变得更具动态交互性，但只有在将其与真实使用场景（例如监控仪表盘）结合时才有意义。用户在多维看板上筛选结果，以借助 Vue 3 的响应式状态管理更迅速地做出响应。在这种模式下，ECharts 能在几毫秒内完成重渲染，并将复杂的统计数据转换为执行趋势折线或科目占比图表，从而更易于观察。这种沉浸式的可视化数据体验将曾经枯燥、分散的数字转化为决策依据，并使其具备直觉指引性。

\section{本章小结}
本章详细阐述基于 Spring Boot、MyBatis-Plus 以及 Vue 3/ECharts 的全栈技术链，采用自底向上的方式，通过概述系统的整体技术架构来展开。其中，Spring Boot 与 MyBatis-Plus 共同提供从底层数据持久化到业务聚合的处理逻辑；而 Vue 3 与 ECharts 则在表现层实现将多维模型动态映射到可视化仪表盘。一方面，这种集成使得“数据收集——存储——分析——可视化”的闭环成为可能；另一方面，它也为后续章节中关于系统设计与实现的讨论奠定了坚实的技术基础。
